\subsection{The actual wild source over its fixed free quotient}

\begin{statement}{E.77}{Proved}{A relative equivariant presentation and its nonlinear obstruction}
Retain the actual local groups, quotient and involution
of E.74--E.75. There are actual generators
$t,t',r,r'$ of $D=G_{\mathbb Q_3(\zeta_3)}(3)$ for which
\[
\begin{gathered}
 ct c=t',\qquad cr c=r',\qquad
 \pi(t)=\tau,\quad\pi(t')=\tau',\quad
 \pi(r)=\pi(r')=1,\\
 D\ \simeq\
 \left\langle t,t',r,r'\ \middle|\
 (trt^{-1}r^{-7})(t'r't'^{-1}r'^{-7})^{-1}=1
 \right\rangle_{\mathrm{pro}\text{-}3}.
\end{gathered}
 \tag{E77.1}
\]
The displayed isomorphism commutes with the fixed
map to $P$ and with $c$. Its construction retains
the actual local relation and a convergent sequence
of changes of generators over $P$.
In particular it does not replace that quotient by
another free quotient with the same rank or signature.

The generators $t,t'$ define a new actual section
$j_1:\Gamma\to\mathcal D$ of the same quotient.
Let $j_0$ be the section fixed in E.75. Its exact
comparison with $j_1$ is the continuous cocycle
\[
 h(g)=j_1(g)j_0(g)^{-1}\in R_3,\qquad
 h(c)=1,\qquad
 h(gg')=h(g)\,j_0(g)h(g')j_0(g)^{-1}.
 \tag{E77.2}
\]
The final $r,r'$ have the same classes in
$R_3^{\rm ab}$ as E.75's $r_*,r_*'$.
All choices in (E77.1)--(E77.2) are made once
in the actual source.

Let $W$ be any of the actual geometric pro-three
quotients in E.41, with its actual $\Gamma$ action,
and write $T(w)=\tau(w)$ and $J(w)=c(w)$.
The full continuous local section set, modulo
geometric conjugacy, has the following exact
compact coordinates:
\[
 \frac{\left\{(u,a)\in W^2:
          J\bigl(uT(a)u^{-1}a^{-7}\bigr)
             =uT(a)u^{-1}a^{-7}\right\}}{W^c}.
 \tag{E77.3}
\]
The residual action is
\[
 w:(u,a)\longmapsto
       \bigl(wuT(w)^{-1},\,waw^{-1}\bigr),
       \qquad w\in W^c.
 \tag{E77.4}
\]
The inverse sends $c$ to the prescribed $c$-lift
and gives the four actual source generators the
section values
\[
 z(t)=u,\quad z(t')=J(u),\quad
 z(r)=a,\quad z(r')=J(a).
 \tag{E77.5}
\]
Equations (E77.1) and (E77.3) prove directly that
these values extend to a continuous section.
The quotient in (E77.3) has its compact orbit
topology; the correspondence and its inverse
are homeomorphisms.

This also evaluates each wild central obstruction
from actual lower section coordinates. For the
central extension $W_r\to W_{r-1}$ with kernel
$A=\mathcal A_r(W)$, choose the retained Hall lifts
$u_0,a_0$ of a pair in (E77.3) at level $r-1$ and put
\[
 v_0=u_0T(a_0)u_0^{-1}a_0^{-7},\qquad
 \omega=v_0J(v_0)^{-1}\in A^-.
 \tag{E77.6}
\]
The local obstruction in the invariant coordinate
of E.74 is
\[
 [\omega/2]\ \in\
 A/\bigl((T-7I)A+(J+I)A\bigr).
 \tag{E77.7}
\]
For the actual lifts $u=u_0b$, $a=a_0d$,
$b,d\in A$, the exact new defect is
\[
 \omega+(I-J)(T-7I)d.
 \tag{E77.8}
\]
Thus the image consists precisely of the lower
pairs whose displayed defect is in the negative
of that explicit image. Each nonempty fibre is
the compact group
$\mathcal C_r(W)\oplus\mathcal K_A$ of E.75,
with its actual correction coordinates.
In particular the source relation in this
criterion is fully specified. Its evaluation
retains the actual geometric operators $T,J$.

\emph{Manuscript: \S51.16, folios733v--736v,
scans172--178, formulas(371)--(379); Part E(c),(d)
and follow-up(4). This supplies the nonlinear
local source and exact obstruction map following
E.74--E.75. It does not replace the remaining
global arithmetic evaluations in E.23.}
\end{statement}

\begin{proof}
\emph{An actual initial relative presentation.}
Start with the generators and minimal relation
module used to prove E.75. Let
\[
 F=\langle t_0,t_0',r_0,r_0'\rangle_{\mathrm{pro}\text{-}3},
 \qquad N=\ker(F\to D),\qquad
 K=\ker(F\to P).
 \tag{E77.9}
\]
The involution swaps the two pairs and the
horizontal generators initially lift $j_0$.
All normal subgroups and commutator subgroups
in this proof are closed. The universal relative
abelian quotient is the actual split group
\[
 F/K'\simeq \Lambda^2\rtimes P,\qquad
 K^{\rm ab}\simeq\Lambda^2,\qquad
 \Lambda=\mathbb Z_3[[P]].
 \tag{E77.10}
\]
Here a row $(\lambda_1,\lambda_2)$ represents
the corresponding completed sum of conjugates
of $r_0,r_0'$. For an explicit proof, assign
$t_0,t_0'$ their marked images in $P$ and assign
$r_0,r_0'$ the two module basis vectors.
The free universal property gives the map in
(E77.10). Its inverse sends a point mass
$\delta_p e_i$ to $p r_i p^{-1}$ modulo $K'$.
The continuous action extends this map to the
completed group ring. The two maps are inverse
on all marked generators and then on their
closed generated groups.

The equivariant matrix used in E.75 is an
actual change of generators in this free
relative presentation. Indeed every prescribed
row of $\Lambda^2$ lifts to an element of $K$.
Lift the first row of the inverse basis matrix
and take its $c$-conjugate for the second row.
Together with $t_0,t_0'$ these elements are a
free basis of $F$: their reduction on the
Frattini quotient is an invertible block
matrix. The resulting endomorphism is onto
by the Frattini criterion and is injective
because a finitely generated profinite group
is Hopfian. One direct proof of the latter
fact uses the finite set of open normal
subgroups of any bounded index: a surjective
endomorphism permutes that set, and hence
is injective on every associated finite
characteristic quotient.
This change fixes the quotient $P$ and
commutes with $c$.

After this change the oriented top generator
$e\in N^{\rm ab}$ has relative Fox row
$(\tau-7,-\tau'+7)$ and satisfies $ce=-e$,
exactly as in E.75. Lift $e/2$ to an actual
$b\in N$ and take
\[
 \rho=b\,c(b)^{-1}.
 \tag{E77.11}
\]
Its class is $e$, it normally generates $N$
by the relative Frattini criterion, and
$c(\rho)=\rho^{-1}$ holds in the free group
itself. A lift of a chosen $r_*$ in E.75 may
be used: any discrepancy of its relative
module row is in the image of $N^{\rm ab}$
and can be corrected by an element of $N$.
Thus the initial actual images of the two
vertical generators can be exactly $r_*,r_*'$.

In this new basis put
\[
 \rho_0=(trt^{-1}r^{-7})
              (t'r't'^{-1}r'^{-7})^{-1}.
 \tag{E77.12}
\]
It too is exactly inverted by $c$. Its row
in $K^{\rm ab}$ is the same oriented row
$(\tau-7,-\tau'+7)$. Consequently
$\rho_0^{-1}\rho\in K'$.

\emph{The integral filtration and its relative kernel.}
Write $F_1=F$ and
$F_{n+1}=F_n^3[F_n,F]$.
We use the commutator $[x,y]=xyx^{-1}y^{-1}$.
The associated graded object is
\[
 \operatorname{gr}F
   =\mathfrak L_{\mathbb F_3[\pi]}
          (T,T',R,R'),\qquad \deg\pi=1.
 \tag{E77.13}
\]
Cubing induces multiplication by $\pi$.
This is also the odd-prime case of the
original lower-central calculation in
\cite[\S2, Propositions1--2, pp.112--113]{A_Labute1967}.
We give the needed integral justification.
Use the Magnus map into the completed
noncommutative power-series algebra over
$\mathbb Z_3$, assigning $3$ and every
variable weight one. Its graded coefficient
ring is $\mathbb F_3[\pi]$.
For $g\in F_n$, the leading term of $g^3$
is three times the leading term of $g$:
the other two terms in $(1+X)^3$ have
strictly higher weight, also for $n=1$.
The commutator leading term is the Lie
bracket. The same formulas show additivity
of the cube map on each layer and
$\pi[x,y]=[\pi x,y]=[x,\pi y]$.
Induction on the defining filtration
therefore spans each layer by the elements
$\pi^a H$ for integral Hall commutators $H$.
Their Magnus leading polynomials are
independent: Hall polynomials of each
ordinary degree have distinct triangular
leading words, and the powers of $\pi$
separate the remaining weights. This
proves (E77.13), with no division by three.
The filtration tends to one, since its
image in every finite three-group does.

Let $\mathfrak J$ be the Lie ideal generated
by $R,R'$ in (E77.13). The actual induced
filtrations satisfy
\[
 \operatorname{gr}K=\mathfrak J,\qquad
 \operatorname{gr}K'=\mathfrak J'.
 \tag{E77.14}
\]
For clarity, the second equality is for
$K'\cap F_n$, not a reindexed intrinsic
filtration on $K'$.
Here is the strict comparison. Put
$\mathfrak m=(3,\tau-1,\tau'-1)\subset\Lambda$.
The relative Fox map $K\to\Lambda^2$ in
(E77.10) is filtered from $K\cap F_n$ to
$\mathfrak m^{n-1}\Lambda^2$. To check this
without an intersection assumption, filter
$\Lambda^2\rtimes P$ by
$\mathfrak m^{n-1}\Lambda^2\rtimes P_n$.
The formula for a commutator and
$(1+p+p^2)v$ for a cube show that this is
a lower three-central filtration.
Its graded relative module is
$\mathbb F_3[\pi]\langle T,T'\rangle^2$.

The ideal $\mathfrak J$ is the free Lie
algebra on the free module with basis
\[
 \operatorname{ad}_{X_1}\cdots
 \operatorname{ad}_{X_d}R_i,\qquad
 X_j\in\{T,T'\},\quad R_i\in\{R,R'\}.
 \tag{E77.15}
\]
An elementary universal-property proof
constructs this free Lie algebra and lets
$T,T'$ act by prepending their adjoints
as derivations. Its semidirect product
with the free horizontal Lie algebra
receives and sends back all four free
generators. The two maps are inverse.
Thus its abelianization is precisely
$\mathbb F_3[\pi]\langle T,T'\rangle^2$,
and the graded relative Fox map has
kernel $\mathfrak J'$.
If $k\in K'\cap F_n$, its Fox value is
zero, so its initial form lies in
$\mathfrak J'$. Conversely every bracket
and every $\pi$ multiple defining a
homogeneous element of $\mathfrak J'$
lifts to a commutator or cube in $K'$.
This proves both inclusions in the
second equality of (E77.14).
The first follows either from the same
argument or from the actual section
$P\to F$, which makes $F_n\to P_n$
surjective with the stated graded kernel.

\emph{The relative correction map.}
Suppose $\rho=\rho_0 e_n$ with
$e_n\in K'\cap F_n$.
At degree two the error is a multiple
of $[R,R']$. The paired shear
\[
 t\longmapsto t(r')^a,\qquad
 t'\longmapsto t'r^a
 \tag{E77.16}
\]
changes the initial form of $\rho_0$ by
$-2a[R,R']$. Choose $a$ modulo three to
remove the error. This is an invertible
triangular change of the Frattini basis,
and it preserves the actual quotient
and the two vertical module classes.
Thereafter $n\ge3$.

Change $t,t'$ by right factors in
$K\cap F_{n-1}$ and $r,r'$ by right
factors in $K'\cap F_{n-1}$.
For their initial forms
$(x,x',y,y')\in
\mathfrak J_{n-1}^2\oplus(\mathfrak J'_{n-1})^2$,
the initial change of $\rho_0$ is
\[
 \mathcal L(x,x',y,y')
  =[x,R]-[x',R']
       +(\operatorname{ad}T-2\pi)y
       -(\operatorname{ad}T'-2\pi)y'.
 \tag{E77.17}
\]
In particular the coefficient is $2\pi$
because $7-1=2\cdot3$.
The ordinary commutator expansion proves
this formula; terms from
$\binom72=21$ have an extra factor of
three and are above the displayed weight.
It is the relative version of the
variation on pp.116--118 of
\cite{A_Labute1967}.

The image of (E77.17) is all of
$\mathfrak J'_n$. This assertion is
integral in the graded coefficient
ring and has a direct proof.
Using Jacobi, every bracket in the
free Lie algebra (E77.15) is a sum
of brackets $[v,b]$ where $v$ is one
of its displayed free generators and
$b\in\mathfrak J$. This follows by
repeatedly expanding
$[[v_1,v_2],b]=[v_1,[v_2,b]]
 -[v_2,[v_1,b]]$; coefficients and
powers of $\pi$ may be placed on $b$.
Modulo the image of (E77.17),
\[
 [\operatorname{ad}T(v),b]
       \equiv 2\pi[v,b]-[v,\operatorname{ad}T(b)],
 \tag{E77.18}
\]
and the same identity holds for $T'$.
Indeed their difference is
$(\operatorname{ad}T-2\pi)[v,b]$,
an allowed term. Induction on the
horizontal word in $v$ reduces its
first entry to $R$ or $R'$.
The resulting bracket is in
$[\mathfrak J_{n-1},R]+
 [\mathfrak J_{n-1},R']$, also in
the image. All operations preserve
the total weight, which proves
surjectivity in every degree.

This calculation preserves the actual
involution. Let $c_\mathrm{dom}$ interchange
the paired correction entries and apply
$c$ to each. Then
$\mathcal Lc_\mathrm{dom}=-c\mathcal L$.
The initial error is anti-invariant:
the exact identities
$c(\rho)=\rho^{-1}$ and
$c(\rho_0)=\rho_0^{-1}$ give this
after passing to the central layer.
Average a solution with its
$c_\mathrm{dom}$ transform, dividing
by two. Lift its first and third
entries to the stated actual closed
subgroups and use their $c$-conjugates
for the other entries. This is an
exactly paired change of generators.
Choose its sign to cancel the error.

All these changes are over $P$.
Horizontal changes by $K$ act on
$K^{\rm ab}$ by the same marked
$P$ action, and vertical changes
by $K'$ do not change its basis.
The two relative rows therefore
continue to agree exactly, so the
next error remains in $K'$ as
well as in $F_{n+1}$.
For $n\ge3$ the changes are
Frattini automorphisms.
Their depths tend to infinity;
the four generators converge to
a basis, because their limiting
Frattini matrix remains invertible.
The limiting relation is exactly
$\rho=\rho_0$ in that basis.
This proves (E77.1) as an isomorphism
over the fixed $P$, with $c$ retained.
It also proves the asserted
$R_3^{\rm ab}$ comparison.

\emph{Change of section and nonlinear coordinates.}
These are the full local section sets for
$G_{\mathbb Q_3}$. Indeed the kernel of its
map to $\mathcal D$ acts trivially on
$W$, and E.74's Kummer comparison proves
that this kernel has no quotient $C_3$.
Any nontrivial continuous homomorphism
from it to a pro-three group would have
a nontrivial finite three-group image,
and hence a quotient $C_3$. Its restriction
under every cocycle is therefore trivial.
All the cocycles in question factor
through the displayed actual local source.

Since $\pi$ maps the new horizontal
pair to the free basis $\tau,\tau'$,
their map from $P$ is injective
and is a section. Together with
the unchanged $c$ it gives $j_1$.
Multiplication proves (E77.2).
For any actual local section write,
as in E.75,
$z(rj_0(g))=\theta(r)u_0(g)$.
The exact new coordinates are
\[
 u_1(g)=\theta(h(g))u_0(g),\qquad
 z(rj_1(g))=\theta(r)u_1(g).
 \tag{E77.19}
\]
Conversely $u_0(g)=
\theta(h(g))^{-1}u_1(g)$.
The equivariance of $\theta$
changes by the inner conjugation
by $\theta(h(g))$, as is seen
by applying $\theta$ to
$j_1(g)rj_1(g)^{-1}$.
These formulas retain the
actual change of splitting.

Coprime conjugacy for $C_2$
normalizes the section on $c$.
For completeness this follows
at finite three-group quotients
by the conjugacy of order-two
complements: induction on the
nontrivial invariant centre
reduces it to an abelian group,
where division by two kills
the cocycle. Compactness
gives compatible conjugators
in their inverse limit.
The residual conjugators are
exactly $W^c$.
The source relation (E77.1),
evaluated in $W\rtimes\Gamma$,
is $vJ(v)^{-1}$, where
$v=uT(a)u^{-1}a^{-7}$.
This proves (E77.3)--(E77.5)
in both directions.
The generators identify the
continuous Hom space with a
closed subspace of $W^4$,
and all conjugating groups
are compact. The obtained
continuous bijections between
compact Hausdorff orbit spaces
are therefore homeomorphisms.

Finally, in a central layer
$\omega$ in (E77.6) lies in
$A^-$ because $J(\omega)=
\omega^{-1}$. Replacing
$a_0$ by $a_0d$ changes $v_0$
by $(T-7I)d$; replacing $u_0$
by $u_0b$ has no effect.
This proves (E77.8) exactly,
including its sign.
The oriented relator was
retained throughout, so the
index-two comparison in E.75
is still $y\mapsto y/2$,
giving (E77.7).
When the defect is removed,
the differences are $b\in A$
and $d\in\mathcal K_A$;
central residual conjugation
by $w\in A^+$ changes $b$
by $(I-T)w$ and leaves $d$
fixed. There is no additional
twisted lower stabilizer by
E.74's positive-weight argument.
The resulting fibre is exactly
$A/(T-I)A^+\oplus\mathcal K_A
=\mathcal C_r(W)\oplus\mathcal K_A$,
with the retained Smith
representatives and inverse
maps of E.41 and E.75.
\end{proof}
